\section{Conclusion}
This study explore the application of YOLO-based object detection for detected and classified Khmer traffic sign.
The experimental results on the Khmer dataset demonstrated both YOLOv8 and YOLOv10 models achieved strong detection
performance, with YOLOv8 showing a small overall advantage across the evaluation measure. YOLOv8 was able to detect
and classify Khmer traffic signs even when the signs had different sizes and the images had different lighting conditions
and image quality.

This study contributes to the development of Khmer traffic sign recognition by evaluating modern object detection models
using a Khmer traffic sign dataset. The results demonstrate the potential of deep learning for automatically detecting and
classifying Khmer traffic signs. Furthermore, this study provides a useful reference for future research and the development
of computer vision-based traffic sign recognition systems for the Cambodian road environment.

For future work, the Khmer traffic sign dataset can be expanded by including more traffic signs images within each classes 
and incorporating additional traffic sign classes to create more comprehensive dataset. Future research could also explore other advanced and lightweight
object detection models to enhance detection performance in real-world conditions including weather conditions, lighting environments,
and road locations. 